% GENERATED FILE. Edit canonical lessons, then run npm run generate:books.
% canonical-source-hash: fnv1a64:de34a842d847643d
% canonical-lessons: KA-C297-jenunona, KA-C297-gubbacci, KA-C297-mosale, KA-C297-cirate, KA-C297-bala

\chapter{Bee, Sparrow, Crocodile, Leopard, Tail}
\label{ch:ka-bee-sparrow-crocodile-leopard-tail}
\hlchaptermodality{297}

\begin{chapteropening}
\textbf{By the end of this chapter:} \emph{I can say a bee, a sparrow, a crocodile, a leopard and a tail.}

Complete the last lesson of chapter 297: I can say a bee, a sparrow, a crocodile, a leopard and a tail.
\end{chapteropening}

\section[\texorpdfstring{\kn{ಜೇನುನೊಣ} (jēnunoṇa)}{ಜೇನುನೊಣ (jēnunoṇa)}]{\kn{ಜೇನುನೊಣ} (jēnunoṇa) — a bee}

\label{lesson:KA-C297-jenunona}



\begin{quote}
Before the new one: say the Kannada for a helmet, then the Kannada for a dictionary.
\end{quote}

\subsection*{You'll want to know: \kn{ಜೇನುನೊಣ}}
\textbf{\kn{ಜೇನುನೊಣ}} — \emph{jēnunoṇa} — \textquotedblleft{}a bee\textquotedblright{}.

\begin{grammarlens}[title={What you've built}]
One more word for this chapter.
\end{grammarlens}

\subsection*{Your turn}
\begin{itemize}
\raggedright
  \item \emph{Say it:} \emph{jēnunoṇa}
  \item \emph{Say it:} \emph{jēnunoṇa}, once more
  \item \emph{Say it:} say \emph{jēnunoṇa}
  \item \emph{Recall:} say the Kannada for a helmet, then the Kannada for a dictionary, then say \emph{jēnunoṇa} again
\end{itemize}

\begin{tcolorbox}[breakable,colback=teal!4,colframe=teal!35!black,title={Before you move on}]
What does \kn{ಜೇನುನೊಣ} mean? (\textquotedblleft{}A bee\textquotedblright{}.) Say it once more.
\end{tcolorbox}
\section[\texorpdfstring{\kn{ಗುಬ್ಬಚ್ಚಿ} (gubbacci)}{ಗುಬ್ಬಚ್ಚಿ (gubbacci)}]{\kn{ಗುಬ್ಬಚ್ಚಿ} (gubbacci) — a sparrow}

\label{lesson:KA-C297-gubbacci}



\begin{quote}
Before the new one: say the Kannada for a dictionary, then the Kannada for a bee.
\end{quote}

\subsection*{You'll want to know: \kn{ಗುಬ್ಬಚ್ಚಿ}}
\textbf{\kn{ಗುಬ್ಬಚ್ಚಿ}} — \emph{gubbacci} — \textquotedblleft{}a sparrow\textquotedblright{}.

\begin{grammarlens}[title={What you've built}]
One more word for this chapter.
\end{grammarlens}

\subsection*{Your turn}
\begin{itemize}
\raggedright
  \item \emph{Say it:} \emph{gubbacci}
  \item \emph{Say it:} \emph{gubbacci}, once more
  \item \emph{Say it:} say \emph{gubbacci}
  \item \emph{Recall:} say the Kannada for a dictionary, then the Kannada for a bee, then say \emph{gubbacci} again
\end{itemize}

\begin{tcolorbox}[breakable,colback=teal!4,colframe=teal!35!black,title={Before you move on}]
What does \kn{ಗುಬ್ಬಚ್ಚಿ} mean? (\textquotedblleft{}A sparrow\textquotedblright{}.) Say it once more.
\end{tcolorbox}
\section[\texorpdfstring{\kn{ಮೊಸಳೆ} (mosaḷe)}{ಮೊಸಳೆ (mosaḷe)}]{\kn{ಮೊಸಳೆ} (mosaḷe) — a crocodile}

\label{lesson:KA-C297-mosale}



\begin{quote}
Before the new one: say the Kannada for a bee, then the Kannada for a sparrow.
\end{quote}

\subsection*{You'll want to know: \kn{ಮೊಸಳೆ}}
\textbf{\kn{ಮೊಸಳೆ}} — \emph{mosaḷe} — \textquotedblleft{}a crocodile\textquotedblright{}.

\begin{grammarlens}[title={What you've built}]
One more word for this chapter.
\end{grammarlens}

\subsection*{Your turn}
\begin{itemize}
\raggedright
  \item \emph{Say it:} \emph{mosaḷe}
  \item \emph{Say it:} \emph{mosaḷe}, once more
  \item \emph{Say it:} say \emph{mosaḷe}
  \item \emph{Recall:} say the Kannada for a bee, then the Kannada for a sparrow, then say \emph{mosaḷe} again
\end{itemize}

\begin{tcolorbox}[breakable,colback=teal!4,colframe=teal!35!black,title={Before you move on}]
What does \kn{ಮೊಸಳೆ} mean? (\textquotedblleft{}A crocodile\textquotedblright{}.) Say it once more.
\end{tcolorbox}
\section[\texorpdfstring{\kn{ಚಿರತೆ} (cirate)}{ಚಿರತೆ (cirate)}]{\kn{ಚಿರತೆ} (cirate) — a leopard}

\label{lesson:KA-C297-cirate}



\begin{quote}
Before the new one: say the Kannada for a sparrow, then the Kannada for a crocodile.
\end{quote}

\subsection*{You'll want to know: \kn{ಚಿರತೆ}}
\textbf{\kn{ಚಿರತೆ}} — \emph{cirate} — \textquotedblleft{}a leopard\textquotedblright{}.

\begin{grammarlens}[title={What you've built}]
One more word for this chapter.
\end{grammarlens}

\subsection*{Your turn}
\begin{itemize}
\raggedright
  \item \emph{Say it:} \emph{cirate}
  \item \emph{Say it:} \emph{cirate}, once more
  \item \emph{Say it:} say \emph{cirate}
  \item \emph{Recall:} say the Kannada for a sparrow, then the Kannada for a crocodile, then say \emph{cirate} again
\end{itemize}

\begin{tcolorbox}[breakable,colback=teal!4,colframe=teal!35!black,title={Before you move on}]
What does \kn{ಚಿರತೆ} mean? (\textquotedblleft{}A leopard\textquotedblright{}.) Say it once more.
\end{tcolorbox}
\section[\texorpdfstring{\kn{ಬಾಲ} (bāla)}{ಬಾಲ (bāla)}]{\kn{ಬಾಲ} (bāla) — a tail}

\label{lesson:KA-C297-bala}



\begin{quote}
Before the new one: say the Kannada for a crocodile, then the Kannada for a leopard.
\end{quote}

\subsection*{You'll want to know: \kn{ಬಾಲ}}
\textbf{\kn{ಬಾಲ}} — \emph{bāla} — \textquotedblleft{}a tail\textquotedblright{}.

\begin{grammarlens}[title={What you've built}]
One more word for this chapter.
\end{grammarlens}

\subsection*{Your turn}
\begin{itemize}
\raggedright
  \item \emph{Say it:} \emph{bāla}
  \item \emph{Say it:} \emph{bāla}, once more
  \item \emph{Say it:} say \emph{bāla}
  \item \emph{Recall:} say the Kannada for a crocodile, then the Kannada for a leopard, then say \emph{bāla} again
\end{itemize}

\begin{tcolorbox}[breakable,colback=teal!4,colframe=teal!35!black,title={Before you move on}]
What does \kn{ಬಾಲ} mean? (\textquotedblleft{}A tail\textquotedblright{}.) Say it once more.
\end{tcolorbox}
